\subsection{场景 B 初解的重新构造}
场景 A 已找到的 $(\pi,m,\sigma)$ 可以作为 B 的一个\emph{结构种子}，却不能把 $T^A(p)$ 搬进 B 的候选排名。变化不只在 $1000$ 与 $500$ 的数字：Task 组成、同核缓存是否保留、跨核 COPY 的生成单位、释放时刻都变了。对每张图、每个多核配置，先在 B 的官方规则下评价移植方案，同时保留整图单 Task 的合法备选；然后用 B 的成本模型重新把分量、共享输入链及深度窗口等候选映射到核心。这样的初始化保留问题一划分可能带来的结构价值，也让 B 专有的同核复用可以重新改变排名。

对弱连通分量整体分核，场景 B 的收益既来自分量间天然无依赖，也来自把同一分量内的相关张量留在一个核。对超过平均 Pipe 份额的主导分量，尝试沿拓扑顺序有限切开，恢复可并行空间。共享输入候选则把同一张量的消费者尽量安排在较少核心，使式\eqref{eq:q2-destination-count}的目的地数下降。长链与深度窗口是另一种结构视角：前者保护可能控制最终输出的连续依赖，后者在近似同层的算子间寻找并行。所有结构都只是\emph{初解来源}，没有一种可单凭静态图保证完成时间更短。

已采用的结构准入依赖一个可复核指标：最大弱连通分量的算子数占比
\begin{equation}\label{eq:q2-dominant-ratio}
\chi(G)=\frac{\max_{C\in\operatorname{WCC}(G)}|C|}{|\mathcal O_c|}.
\end{equation}
当 $\chi(G)>0.5$ 时，允许对候选菜单末段采用针对主导分量的另一种有界布局；当其不超过该阈值，保持原基础菜单。阈值并非物理常数，也没有被解释为通用最优切分点；它是全量实验前固定的结构门控规则。更关键的是，前部已有效的初解机会保持不变，新增布局只使用原菜单中预留的后部机会。若某个布局与移植种子、整图方案或前部布局完全相同，依据完整方案指纹识别重复，并用一个结构不同的旧候选补回该机会；没有把重复评价算作新的发现，也没有增加总预算。

\subsection{联合映射和有界普通邻域}
在相同划分下改变一个子图的核心，会同时改变跨核目的地集合、同核 Task 内的张量生命周期、跨核 COPY 释放以及相邻子图的资源排序；因此映射不能在划分之后靠“每核任务数相同”一次性决定。构造初解时，对就绪子图 $r$ 和候选核心 $c$，用场景 B 的释放估计 $\widehat R_B(r,c)$ 与核内负载/驻留估计 $\widehat d_B(r,c)$ 形成
\begin{equation}\label{eq:q2-eft}
\widehat F_B(r,c)=\widehat R_B(r,c)+\widehat d_B(r,c),
\qquad
\widehat R_B(r,c)=
\max\left\{\widehat F_{\rm same\ core},
\max_{a\in\operatorname{Pred}(r)}
\bigl(\widehat F_a+\widehat c_{ar}(c)\bigr)\right\}.
\end{equation}
这里 $\widehat c_{ar}$ 对同核数据采用驻留代理，对跨核数据采用按消费核心去重后的传输和同步代理。它不把 $500$ 周期加到所有前驱上；同一时间可并行完成的等待取最大。$\widehat F_B$ 的用途是按最早预计完成的方向生成候选；正式 B 评价会重新展开 COPY、Pipe 和物理缓存。

基础结构评价结束后，求解保留至多四个普通联合邻域机会。对当前实测较优方案，在一个小关联区域中选子图 $r$、目标核心 $c$ 及合法插入位置 $g$，形成
\begin{equation}\label{eq:q2-neighborhood}
\mathcal N_J(p)=
\left\{
\operatorname{RepairOrder}\bigl(\operatorname{MoveCore}(p,r,c),g\bigr)
:r\in\mathcal R(p),\ c\ne m(r),\ g\in\mathcal G_{r,c}
\right\}.
\end{equation}
$\mathcal R(p)$ 是受预算和结构规则限制的少数相关子图，$\mathcal G_{r,c}$ 是不会破坏核上顺序与数据依赖的有限位置集合。$\operatorname{RepairOrder}$ 必须在迁核后同步修复被影响任务的局部顺序，否则孤立迁核可能只增加等待而看不到联合动作的价值。实现从这个邻域提出有限不同候选，逐个做覆盖和扩展商图无环检查，再由官方 B 仿真裁决。式\eqref{eq:q2-neighborhood}定义了可理解的邻域形状，不表示对其中所有组合穷举。

算法在每个多核配置保留十二个候选评价位置：移植方案和整图保底、六个结构候选以及四个有界普通联合动作。重复方案可由带输入、配置、评价代码身份的缓存直接复用相同官方结果，但它在台账中仍占原位置，只有明确定义的重复候选补位才改变菜单的有效覆盖数；不能因缓存命中就在对照组偷偷增加评价机会。最终胜出方案再由原始 B 评价过程独立回放，检查完整结果而不仅比较 Makespan。这个固定流程使研究中临时尝试的关键链邻域、替代通信代理或生命周期排序不被误写成所交成绩的组成部分。

\subsection{搜索正确性和代价控制}
只要移植或整图保底中至少有一个合法方案，逐项过滤并以官方结果更新 incumbent 的算法返回合法方案。其形式化性质是：设 $\mathcal A_B$ 为本次实际完成官方 B 评价的合法候选集合，返回 $p_B^\star$ 满足
\begin{equation}\label{eq:q2-finite-optimality}
p_B^\star\in\mathcal A_B,\qquad
\operatorname{lex}\bigl(T^B(p_B^\star),D^B(p_B^\star)\bigr)
=\min_{p\in\mathcal A_B}
\operatorname{lex}\bigl(T^B(p),D^B(p)\bigr).
\end{equation}
证明直接由每次把新候选与已有最优实测方案作字典序比较的归纳过程得到。该性质仅限\emph{实际评价的有限候选集}，不能扩展为 $\mathcal P_B(k)$ 的全局最优。若代理把真正优秀的候选排在第十三位，它仍会被预算截断；这正是选择后悔与预算机会成本必须审计的原因。

结构分析、拓扑排序和指纹计算各自随图规模多项式增长；关键成本仍在官方核内准备、物理内存检查与共享 DDR 事件模拟。十二个位置的上限约束的是正式候选机会，而不是保证所有图在固定墙钟秒数内结束。已有成效的基础候选先保留，新邻域若只造成菜单替换而没有形成不同方案，扩大的搜索种类反而可能使结果变差。论文因此只陈述当前经过全百图核验的统一规则，不以开发阶段某几张图的局部改善作为泛化成绩。
